% ============================================================
% Chapter 3: System Analysis and Design
% ============================================================

\chapter{System Analysis and Design}

\section{System Analysis}

\subsection{Requirement Analysis}

\subsubsection{Functional Requirements}

The functional requirements define what the system should do. They are represented using Use Case diagrams and descriptions.

\paragraph{Use Case Diagram}

Figure \ref{fig:use-case} shows the main use cases of the Personal Finance Manager system.

\begin{figure}[H]
    \centering
    % Placeholder for use case diagram
    \includegraphics[width=\textwidth]{diagrams/use-case-diagram.png}
    \caption{Use Case Diagram for Personal Finance Manager}
    \label{fig:use-case}
\end{figure}

\paragraph{Use Case Descriptions}

\begin{table}[H]
\centering
\caption{Use Case: Add Expense via Natural Language}
\label{tab:uc-add-expense}
\begin{tabular}{|p{4cm}|p{9cm}|}
\hline
\textbf{Use Case ID} & UC-001 \\
\hline
\textbf{Use Case Name} & Add Expense via Natural Language \\
\hline
\textbf{Actor} & User \\
\hline
\textbf{Description} & User adds an expense by typing in natural language \\
\hline
\textbf{Precondition} & User is logged in and has selected a group \\
\hline
\textbf{Main Flow} & 
1. User types expense in natural language (e.g., ``biryani 500'')\\
2. System parses the input using NLP\\
3. System extracts amount, item, category, and person (if any)\\
4. System displays parsed expense confirmation\\
5. User confirms the expense\\
6. System saves the expense to database \\
\hline
\textbf{Postcondition} & Expense is recorded in the database \\
\hline
\textbf{Alternative Flow} & If parsing fails, system prompts user to rephrase \\
\hline
\end{tabular}
\end{table}

\begin{table}[H]
\centering
\caption{Use Case: Query Expenses}
\label{tab:uc-query}
\begin{tabular}{|p{4cm}|p{9cm}|}
\hline
\textbf{Use Case ID} & UC-002 \\
\hline
\textbf{Use Case Name} & Query Expenses via Natural Language \\
\hline
\textbf{Actor} & User \\
\hline
\textbf{Description} & User queries expenses using natural language \\
\hline
\textbf{Precondition} & User is logged in with expense data available \\
\hline
\textbf{Main Flow} & 
1. User types query (e.g., ``How much on food this month?'')\\
2. System analyzes query intent\\
3. System retrieves relevant expenses\\
4. System generates response using RAG\\
5. System displays response to user \\
\hline
\textbf{Postcondition} & User receives answer to query \\
\hline
\end{tabular}
\end{table}

\subsubsection{Non-Functional Requirements}

\begin{table}[H]
\centering
\caption{Non-Functional Requirements}
\label{tab:nfr}
\begin{tabular}{|p{3cm}|p{10cm}|}
\hline
\textbf{Requirement} & \textbf{Description} \\
\hline
Performance & System should parse expenses within 2 seconds \\
\hline
Availability & System should be available 99\% of the time \\
\hline
Security & User data must be encrypted; authentication required \\
\hline
Usability & Interface should be intuitive and mobile-friendly \\
\hline
Scalability & System should handle 1000+ concurrent users \\
\hline
Responsiveness & UI should adapt to different screen sizes \\
\hline
\end{tabular}
\end{table}

\subsection{Feasibility Analysis}

\subsubsection{Technical Feasibility}

The project is technically feasible because:
\begin{itemize}
    \item React and FastAPI are mature, well-documented technologies
    \item Google Gemini AI provides reliable NLP capabilities via API
    \item Supabase offers a complete backend-as-a-service solution
    \item Capacitor enables cross-platform deployment from a single codebase
\end{itemize}

\subsubsection{Operational Feasibility}

The system is operationally feasible:
\begin{itemize}
    \item Natural language input reduces learning curve
    \item Automatic categorization simplifies user workflow
    \item Mobile-first design aligns with user preferences
\end{itemize}

\subsubsection{Economic Feasibility}

The project is economically feasible:
\begin{itemize}
    \item All core technologies are open-source or have free tiers
    \item Supabase free tier supports up to 500MB database
    \item Gemini API offers free tier for development
    \item Hosting costs are minimal with serverless deployment
\end{itemize}

\subsubsection{Schedule Feasibility}

The project timeline of one semester (4 months) is feasible given the scope and available resources.

\subsection{Object Modeling}

\subsubsection{Class Diagram}

Figure \ref{fig:class-diagram} shows the main classes and their relationships.

\begin{figure}[H]
    \centering
    \includegraphics[width=\textwidth]{diagrams/class-diagram.png}
    \caption{Class Diagram for Personal Finance Manager}
    \label{fig:class-diagram}
\end{figure}

\paragraph{Main Classes:}

\begin{itemize}
    \item \textbf{User}: Represents authenticated users
    \item \textbf{Expense}: Represents individual expense records
    \item \textbf{Group}: Represents expense groups (personal, family, etc.)
    \item \textbf{NLPService}: Handles natural language parsing
    \item \textbf{RAGService}: Handles query processing
    \item \textbf{ExpenseAnalyzer}: Provides expense analysis
\end{itemize}

\subsection{Dynamic Modeling}

\subsubsection{Sequence Diagram: Expense Parsing}

Figure \ref{fig:sequence-parse} shows the sequence of operations for parsing an expense.

\begin{figure}[H]
    \centering
    \includegraphics[width=\textwidth]{diagrams/sequence-diagram.png}
    \caption{Sequence Diagram for Expense Parsing}
    \label{fig:sequence-parse}
\end{figure}

\subsubsection{State Diagram: Expense Transaction}

Figure \ref{fig:state-diagram} shows the states of an expense transaction.

\begin{figure}[H]
    \centering
    \includegraphics[width=\textwidth]{diagrams/state-diagram.png}
    \caption{State Diagram for Expense Transaction}
    \label{fig:state-diagram}
\end{figure}

\subsection{Process Modeling}

\subsubsection{Activity Diagram: Add Expense}

Figure \ref{fig:activity-add} shows the activity flow for adding an expense.

\begin{figure}[H]
    \centering
    \includegraphics[width=0.8\textwidth]{diagrams/activity-diagram.png}
    \caption{Activity Diagram for Adding Expense}
    \label{fig:activity-add}
\end{figure}

\section{System Design}

\subsection{System Architecture}

Figure \ref{fig:system-arch} illustrates the high-level architecture of the Personal Finance Manager system, showcasing the interaction between the client (React/Capacitor), backend (FastAPI), and external services (Gemini, Supabase).

\begin{figure}[H]
    \centering
    \includegraphics[width=\textwidth]{diagrams/system-architecture.png}
    \caption{System Architecture Diagram}
    \label{fig:system-arch}
\end{figure}

\subsection{Refined UML Diagrams}

\subsubsection{Component Diagram}

Figure \ref{fig:component-diagram} shows the major components of the system.

\begin{figure}[H]
    \centering
    \includegraphics[width=\textwidth]{diagrams/component-diagram.png}
    \caption{Component Diagram}
    \label{fig:component-diagram}
\end{figure}

\textbf{Components:}

\begin{itemize}
    \item \textbf{Frontend Components}: React UI components (Chat, Table, Analytics, etc.)
    \item \textbf{API Layer}: FastAPI REST endpoints
    \item \textbf{Service Layer}: NLPService, RAGService, ExpenseAnalyzer
    \item \textbf{External Services}: Gemini AI, Supabase
\end{itemize}

\subsubsection{Deployment Diagram}

Figure \ref{fig:deployment-diagram} shows the deployment architecture.

\begin{figure}[H]
    \centering
    \includegraphics[width=\textwidth]{diagrams/deployment-diagram.png}
    \caption{Deployment Diagram}
    \label{fig:deployment-diagram}
\end{figure}

\textbf{Deployment Components:}

\begin{itemize}
    \item \textbf{Client Devices}: Web browsers, Android devices
    \item \textbf{Frontend Server}: Vercel (React app hosting)
    \item \textbf{Backend Server}: Python/FastAPI on cloud server
    \item \textbf{Database}: Supabase PostgreSQL
    \item \textbf{AI Service}: Google Gemini AI (external)
\end{itemize}

\section{Algorithm Details}

\subsection{NLP Parsing Algorithm}

The NLP parsing algorithm uses a hybrid approach combining regex pattern matching with AI-powered parsing.

\begin{lstlisting}[language=Python, caption=NLP Parsing Algorithm]
def parse_expense(text):
    # Step 1: Preprocess text
    text = preprocess(text)  # Handle "lakh", "crore", "k"
    
    # Step 2: Try regex pattern matching
    for pattern in PATTERNS:
        match = pattern.match(text)
        if match:
            return extract_from_match(match)
    
    # Step 3: Fall back to AI parsing
    ai_result = ai_enhanced_parse(text)
    
    # Step 4: Categorize
    category = detect_category(ai_result.item)
    
    # Step 5: Detect person
    person = detect_person(text, ai_result)
    
    return Expense(
        amount=ai_result.amount,
        item=ai_result.item,
        category=category,
        paid_by=person
    )
\end{lstlisting}

\subsection{Category Detection Algorithm}

Categories are detected using keyword matching:

\begin{lstlisting}[language=Python, caption=Category Detection Algorithm]
CATEGORIES = {
    'Food': ['rice', 'momo', 'biryani', 'lunch', 'dinner'],
    'Transport': ['petrol', 'taxi', 'uber', 'bus', 'train'],
    'Shopping': ['shirt', 'shoes', 'clothes', 'bag'],
    # ... 20+ categories
}

def detect_category(item):
    item_lower = item.lower()
    for category, keywords in CATEGORIES.items():
        for keyword in keywords:
            if keyword in item_lower:
                return category
    return 'Other'
\end{lstlisting}

\subsection{RAG Query Processing}

\begin{lstlisting}[language=Python, caption=RAG Query Processing]
def process_query(query, expenses):
    # Step 1: Analyze query intent
    intent = analyze_intent(query)
    
    # Step 2: Retrieve relevant expenses
    if intent == 'category_query':
        relevant = filter_by_category(query, expenses)
    elif intent == 'time_query':
        relevant = filter_by_time(query, expenses)
    else:
        relevant = expenses
    
    # Step 3: Build context
    context = prepare_context(relevant, query)
    
    # Step 4: Generate response with LLM
    response = generate_response(query, context)
    
    return response
\end{lstlisting}
